A trusted research environment (TRE) that admits a local language model as an analyst needs evidence that the model does not mislead on purpose, not only that its honest errors are rare. A companion study measured honest errors on \N{harvest.tasks.total} synthetic analysis tasks and recommended an analyst configuration. Here we ask what else can be learned about a candidate before admission, in four experiments on Qwen3.5 27B and gpt-oss 20B with stop rules fixed in advance. First, linear probes on a model's internal state told wrong write-ups from right ones better than the model's own verdict on unseen tasks (AUC \N{probe.q27.auc} against \N{probe.q27.verbal} for the 27B), but a probe trained on labels shuffled within tasks did almost as well: the probes had learned which tasks are risky, and within tasks they added nothing. Second, a hidden instruction to bias the write-up without changing any number raised wrong conclusions on unseen tasks from \N{steer.v2.heldout.honest} to \N{steer.v2.heldout.steered} of \N{steer.n.heldout}. A probe and a word-count classifier both told biased write-ups from their honest pairs (AUC \N{steer.v2.probe} and \N{steer.v2.text} within tasks), so the internal state gave no advantage over the text, while checkers that judge correctness caught little of the bias. Third, the recommended configuration did not bias when allowed to think first, and a detector trained on another configuration of the same model flagged \N{slant.3style.flagged.k} of its \N{slant.flagged.n} ordinary write-ups. Fourth, a request for brevity raised a detector's false alarms to \N{style.concise.current.pct} of honest write-ups, and training on three styles restored them. An engineering pilot showed that the model can be fine-tuned, reloaded and restored in containment. These tests characterize behavior under a specified instruction, and do not show that a model trained to deceive, or one with a hidden trigger, would be detected.
